\chapter[Sermon 34. The Effect of Corrupt Doctrine Is Expounded, and That the Angels Let Not the Wind Blow.]{The Effect of Corrupt Doctrine Is Expounded, and That the Angels Let Not the Wind Blow.}
\label{sermon:034}
\markright{Sermon 34. The Effect of Corrupt Doctrine Is Expounded, and That ...}

And all mountains and islands were moved out of their places. And the kings of the earth, and the great men, and the rich men, and the chief captains, and the mighty men, and every slave and every free person hid themselves in caves and in the rocks of the hills, and said to the hills and rocks, “Fall upon us and hide us from the presence of him who sits on the throne, and from the wrath of the Lamb,” for the great day of his wrath has come. And who can endure it?

And after this, I saw four angels standing on the four corners of the earth, holding back the four winds of the earth, so that the winds would not blow on the earth, or on the sea, or on any tree.

\index{Mahomet}\index{Papacy}Now follows the effect of the corrupt doctrine in men. And hills and islands are moved out of their place; wherein is also a respect had to the earthquake, as though by the earthquake they were removed from their place. And mountains and islands do betoken realms, nations, and people, so steadfast in faith, that as mountains and islands are immovable, and are not shaken by the storms of the sea, so these might seem to be immutable. Nevertheless, at the alteration and corrupting of doctrine, they are now also removed out of their place, and quite overthrown. And such as read histories shall find everywhere that such have been deceived by the crafts of heretics, by the power of Muhammad, and by the hypocrisy of the pope, whom you would not have thought should have been abused, in so much that whole cities and realms have clean revolted. For seducing is of efficacy namely in such as now begin to slip and slide from the rock of the church.

And those who, being unsettled, are removed from the firm foundation, retreat into caves and rocks of hills. For it is impossible for one who does not hold Christ with steadfast faith to find rest. For like a raging sea, he is tossed to and fro. For since he does not have a sure and certain way of life, nor commits himself to be ruled only by the Scriptures, so that he might hold the certainty, he is content to be led by anyone he encounters. Therefore, we see those to whom Christ alone is not all seeking salvation in pilgrimages, in hermitages, in monasticism, in stricter discipline, in satisfactions, and I know not what other follies, or rather blasphemies. And these are truly said to hide themselves in dens and caves of stone, and think they may lie hidden safely within them, make satisfaction for their sins, and please the Lord.

But in describing many kinds of men, he encompasses all states in the world. For of all sorts of men, there have been found not a few, nor of low status, who have taken upon themselves the hermitic and monastic life, and have bound themselves to a strict kind of living. Here, therefore, are reckoned up kings, judges, great men or princes, rich men, merchants, captains over thousands or chieftains, valiant men in this world, bond servants, and free men, whom we call at this day gentlemen. But how many kings and princes and noble gentlemen are depicted in the churches of abbeys, painted in pictures and hung on trees, who have lived some time a monastic life?

\index{John, Saint}But entering monasteries, woods, and wilderness, and adopting a stricter life, along with various satisfactions, pilgrimages, and similar disciplines, have not yet led to a quietness of mind; indeed, they are now more afraid than they were before, and have fallen into utter despair. For in these things wherein they sought quietness, they have found none; no, besides Christ, there is found no quiet nor rest. Those who live in such strictness under the unhappy papistry understand this very well. And the words which John the Apostle recites here are those of people in the greatest distress, even to desperation, where they cry out, “Let the hills fall upon us,” and so forth. For so this word is also used by Hosea in chapter 10 and by Luke in chapter 23. This signifies a conscience most afflicted and most entangled, finding nowhere any comfort or consolation, but desiring nothing other than present destruction, so as to be delivered from the present evil and intolerable grief of mind. Such a thing is that of Turnus with Virgil in the tenth book of the Aeneid.

Moreover, the causes of this fear, despair, and hiding are the appearance of the one who sits on the throne, the wrath of the Lamb, and because they perceive that they cannot stand before God in the day of wrath and God’s vengeance. Therefore, they flee from the face of God, they flee from the Lamb, that they might avoid the vengeance, if they could escape it. The fear of God is commended to us in the scriptures, and those who do not fear God are condemned; but the scripture speaks of a fear joined with true faith and love. For John says, “Love casts out fear.” Even so, the same scripture preaches to us God as just, and shows him to be angry with sin; nevertheless, it declares him to be benign and merciful to those who acknowledge their sins and ask for forgiveness, that his only begotten Son is given by God to humankind, by whose mediation we may come to the throne of God, which otherwise no man may attain to. It preaches Christ, the Son of God, to be the Lamb, that is a propitiation for the sins of the whole world; and that the same calls all to him, excludes no one, but promises and offers to all, all things of life and salvation. But whereas corrupt preachers, friars, and popish priests have forsaken this simple and most pure doctrine, wholesome and full of consolation, and show that God is like Rhadamanthus, a judge inexorable, and set forth Christ rather as one angry than favorable, they do alienate, without doubt, the minds of men from God; that now they may say expressly, “Who is worthy to come into the sight of God?” No man shall be saved before this God most severe, and his Son a judge most righteous. They turn them therefore to sundry means of salvation, they choose mediators and intercessors by whose mediation and merits they may redeem to themselves the favor of the angry deity. But those who rely on anything other than God’s only mediation and intercession of the Son are disappointed of their purpose, and at the length fall into that same desperation. When they perceive that the monastic life and their merits cannot stand before God, they flee from the face of God; and tormented with the pricks of their conscience, they know not what they may do, whither they may turn, where is the true salvation. Therefore, we judge those who through Christ acknowledge the Father as a Father, and through Christ have access to the Father, as favoring them and loving them, to be rightly most blessed; acknowledging truly in the fear of God their sins, but yet with a true faith hoping for remission of sins, knowing that they are through Christ reconciled to God the Father. The monastic, heretical, expiatory, and Pharisaical faction fully acknowledges this doctrine, therefore they are tormented with sorrows that cannot be uttered. If I speak not here of the monasteries or monks of this our time, in whom we see almost no conscience at all, nor other intent than to be addicted to idleness, voluptuousness, and to bear rule. In times past were found men full of conscience, entering into cells and woods, for no other cause than that they might so be saved. Of such spoke the Lord in the gospel: “When they shall say,” he says, “Christ is in the wilderness, go not forth,” and so forth. And I doubt not, but that some simple also at this day for this intent take upon them the monastic life; but they shall find also, as John said here they should prove and try by experience.

Furthermore, this passage might seem to call for an explanation of the signs preceding the final judgment, and of the terror of the wicked, as the Lord preached in a similar manner in Luke 21. But the final judgment will be discussed more fully in chapters 11 and 19 of this book, and elsewhere. And while I do not condemn such an interpretation, it seems to me that the general destinies of the church are here presented together, in which, where corrupt doctrine occupies not the last place, nothing should be spoken of it in general; many things will be addressed in particular in chapter 8 and those that follow, unless this present passage is explained in the same way as it is. Moreover, those things that follow will be better connected, which will not be part of the final judgment, as the matter itself will demonstrate.

And the things that follow in the 7th chapter pertain to the explanation of the sixth seal, or to the discussion of it. And it chiefly recounts three things: how the angels restrain the winds so they do not blow; an innumerable company to be sealed in the midst of corrupt doctrine, which should not perish. And what the state is of those who have departed out of this world either by martyrdom, or else being either undefiled by the corruption so full of enormity, or delivered and purged from the same; which are added for consolation. For this book of Revelation is wonderfully evangelical, most full not only of prophecies, but also of admonitions, exhortations, and most comforting consolations.

First, it is to be explained what is spoken of the restraint of the winds by the angels, so that they should not blow. Wind, as also appears in the scriptures, is used in both a good and an evil sense. For wind is called vain and false doctrine, and a hope conceived of erroneous doctrine, as in Hosea 12, 5, and 22 of Jeremiah. So leaven is called the Pharisaic doctrine, and the hypocrisy springing from it. Paul in 4 Ephesians forbids that we be carried about with every wind of doctrine. And the Holy Spirit is symbolized by wind in chapter 3 of John. And in chapter 2 of Acts, wind is subtle, it pierces, is felt, and is not seen; great is its force, it cools, it dries, gathers clouds, which bring rain and make the earth fertile. Therefore, rightly, by wind is signified the Spirit of God, and the sound doctrine which is of the Spirit of God. Therefore, it is one wind, the Spirit of God which inspires; and there are four winds, that is to say, many from the corners of heaven and parts of the earth, that is to say, preachers dispersed throughout the whole world. Therefore, the doctrine of the Gospel, inspired from all parts of the world, blows or is preached, so that there are many winds, yet all proceeding from one. For there is one and the same Spirit which speaks by the ministers, and gives them sundry graces, 1 Corinthians 12. Briefly, by the blast of winds we understand the free preaching taken from the holy scriptures.

Secondly, we must know that both good and evil angels are mentioned in the Scriptures. Angels, as appeared before, are called ministers. And there are good and evil ministers: the good inspired by God and the good angel; and the evil by the evil angel. And the enemy of the truth stirs up men in all places of the world, in the courts of kings, in the places of judgment, in schools, in colleges, in cities, towns, and villages, which may hinder the free course of God’s word. Therefore, the proclamations of kings and bishops fly to and fro, are proclaimed and set up, prohibiting the reading of the Bible, the preaching of the Gospel, and so forth. And to have some pretense for their evil doing, they forge that the Bible is corrupt in a thousand places, and that heresy is learned and taught out of the same. Therefore also they prohibit and condemn the Bible and the books of the Gospel, the unworthiness of which thing cannot be sufficiently spoken before the church. They do the same that in times past Antiochus, Epiphanes, Diocletian, and other men of the same sort are recorded to have done. The expositors of the Bible in times past deserved exceeding great praise; neither was there any faithful who said the holy book to be corrupted, for that all translations did not agree among themselves. We live therefore at this day in a time most corrupt and most unthankful.

The suppression of reading holy Scripture is the foundation of corrupt doctrine, and of confusing the conscience, and of the despair that follows from it. And by the Earth he understands those men dwelling on Earth; by the Sea and Isles, those men of islands, and those who dwell on the Sea; by trees, men, everywhere shadowed in Scripture by trees. For unless the winds blow, the trees do not flourish, nor does the earth grow green. The Prophet says, “Send forth your spirit, and they shall be created, and you shall renew the face of the earth.” And except the word of God be preached, the minds of men do not grow green, nor are the fruits of good works brought forth by men. And therefore, the angels prohibiting the wind are said to do harm; for in truth, there is nothing more pestilent nor pernicious than the suppressing of the free preaching of God’s word. The Lord by his spirit renew all parts of the world. Amen.
